\clearpage
\section{Institutional Context}
\label{sec:institutional_context}

This chapter establishes the institutional, legal, and sociological context necessary to interpret the labor market outcomes of higher education policies in Brazil. It transitions from a structural overview of the university system to a detailed analysis of the specific policy mechanisms—specifically the \acrshort{prouni} and the Quota Law—and the legal battles that defined them.

\subsection{Higher Education in Brazil: A Structural Duality}

Brazil's university system is among the most expansive in Latin America, characterized by a distinct public-private duality. According to the latest Census of Higher Education, the system comprises 2,580 institutions, of which more than 85\% are private \cite{INEP2024}. While the federal and state public universities are tuition-free and historically elite, the private sector is highly heterogeneous, absorbing the vast majority of the demand for mass higher education.

\begin{figure}[ht]
    \centering
    \includegraphics[width=0.6\linewidth]{Figures/mapa-riqueza.png}
    \caption{\textbf{The Geography of Inequality:} Average declared income by municipality (2020). The spatial concentration of wealth in the South/Southeast mirrors the historical concentration of university seats. Source: FGV Social (2020)}
    \label{fig:income-map}
\end{figure}

\subsection{Historical and Sociological Foundations}

To understand the design of these policies, one must examine the specific sociology of race and the fiscal structure of education that governed 20th-century Brazil.

\subsubsection{The "Inverted Pyramid" of Public Financing}
The necessity of a private-sector solution (ProUni) is rooted in the "inverted pyramid" of public financing \cite{WorldBank2002}. The state allocates a disproportionate share of its budget to the tertiary level—which is free and dominated by the elite—while basic education remains underfunded. In the early 2000s, public expenditure per student in higher education was approximately 13 times higher than in basic education \cite{McCowan2007}. This regressive fiscal structure constrained the expansion of the federal university system, necessitating a policy design like ProUni that leveraged the private sector's marginal cost structure rather than exclusively expanding the capital-intensive public plant.

\subsubsection{The Sociology of Race: From "Democracy" to Structural Racism}
Unlike the United States, Brazil developed a unique ethnical order based on the ideology of "Racial Democracy", which held that extensive procreation between different ethnicities had created a fluid social structure free from rigid ethnical hatreds. However, as \citeA{Telles2004} documents, this ideology masked deep structural racism where economic inequality remained highly stratified.

A critical institutional challenge for policy is the complexity of racial classification. Brazilian classification is based on phenotype rather than ancestry. The \acrshort{ibge} uses five categories: \textit{branca}, \textit{preta}, \textit{parda}, \textit{amarela}, and \textit{indígena}. However, the phenomenon of "money whitening"—where individuals of higher socioeconomic status are more likely to perceive themselves as white—complicates the data \cite{Telles2004}. This structural reality justifies the aggregation of \textit{Preta} and \textit{Parda} into a single "non-white" category for the econometric analysis, as these groups share the structural disadvantage of the economically weaker sector.

\subsection{The ProUni Program: Design and Formalia}
\label{sec:prouni_design}

The \acrfull{prouni} was established by Law No. 11.096 in January 2005 to democratize access to higher education. By leveraging tax incentives, the program facilitated a massive expansion of the private sector's capacity.

As illustrated in Figure \ref{fig:national_expansion}, the aggregate supply of scholarships grew exponentially following its implementation. Between 2005 and 2014, the annual intake of new beneficiaries increased from approximately 100,000 to over 300,000, accumulating a stock of human capital that represents the primary supply shock analyzed in this thesis.

\begin{figure}[ht]
    \centering
    \includegraphics[width=0.9\textwidth]{Figures/fig_2_1_national.png}
    \caption{\textbf{Aggregate Expansion of ProUni Scholarships (2005--2019)}}
    \label{fig:national_expansion}
    \begin{minipage}{0.9\textwidth}
        \vspace{0.1cm}
        \footnotesize{\textit{Note: The figure displays the total number of new scholarships (full and partial) granted annually. The rapid expansion phase (2005--2014) constitutes the primary supply shock analyzed in this thesis. Source: Author's calculation using \acrshort{mec} microdata.}}
    \end{minipage}
\end{figure}

\subsubsection{Fiscal Mechanism: Tax Expenditures}
ProUni is not a direct scholarship payment but a tax expenditure policy. Participating \acp{hei} receive exemptions from four federal corporate taxes—\textbf{IRPJ}, \textbf{CSLL}, \textbf{PIS/Pasep}, and \textbf{COFINS}—in exchange for offering scholarships \cite{Campos2019}. To prevent the proliferation of "diploma mills," the state linked these exemptions to quality regulation via the \acrfull{sinaes}. Institutions that consistently underperform on the National Student Performance Exam (\acrshort{enade}) face suspension from the program.

\subsubsection{Eligibility and Selection Rules}
The program uses a strict means-test and meritocratic selection based on the \acrfull{enem}. To be eligible, a candidate must meet three criteria:
\begin{enumerate}
    \item \textbf{Secondary Education:} Completion of high school exclusively in public schools or as a full scholarship student in private schools.
    \item \textbf{Academic Merit:} A minimum average score of \textbf{450} on the ENEM and a non-zero score on the essay.
    \item \textbf{Income Thresholds:}
    \begin{itemize}
        \item \textbf{Full Scholarship (100\%):} Gross monthly household income per capita of up to \textbf{1.5 minimum wages}.
        \item \textbf{Partial Scholarship (50\%):} Gross monthly household income per capita of up to \textbf{3 minimum wages} \cite{MEC2024}.
    \end{itemize}
\end{enumerate}

\subsubsection{Affirmative Action Mechanisms}
Crucially, ProUni mandates that institutions reserve a percentage of scholarships for self-declared \textit{Pretos}, \textit{Pardos}, and \textit{Indígenas}, proportional to their share in the respective state's population according to the latest Census \cite{MEC2024}. This effectively created a voucher system for the non-white population to access the private sector. Furthermore, while there are no \textit{de jure} gender quotas, the program functions as a \textbf{de facto gender policy}: women constitute the majority of beneficiaries, driven by higher secondary completion rates and occupational sorting into Health and Education courses \cite{BusteloEtAl2021}.

\subsection{Concurrent Policies: FIES and the Quota Law}

The evaluation of ProUni is complicated by two concurrent shocks that altered the supply of higher education. Figure \ref{fig:timeline} illustrates the temporal overlap of these interventions.

\begin{figure}[ht]
    \centering
    \includegraphics[width=1.0\textwidth]{Figures/fig_3_2_timeline.png}
    \caption{\textbf{Timeline of Higher Education Policies}}
    \label{fig:timeline}
    \begin{minipage}{0.9\textwidth}
        \vspace{0.1cm}
        \footnotesize{\textit{Note: The timeline highlights the three major interventions in the Brazilian higher education market. While FIES and the Quota Law overlap with the ProUni sample period, the identification strategy (Chapter 4) isolates the ProUni effect via the staggered adoption of scholarship density.}}
    \end{minipage}
\end{figure}

\subsubsection{FIES: The Credit Expansion}
Parallel to ProUni, the Student Financing Fund (\acrshort{fies}) provided subsidized loans. Between 2010 and 2014, FIES underwent an aggressive expansion with negative real interest rates, fueling a boom in private enrollments \cite{Campos2019}. Unlike ProUni, FIES students graduate with debt liabilities, theoretically altering their reservation wages.

\subsubsection{The Federal Quota Law (Law 12.711/2012)}
In 2012, the state intervened in the public sector with the \acrfull{lei_cotas}. This law mandated that all federal universities reserve 50\% of admissions for public school graduates, with sub-quotas for income and race \cite{RiehlMachadoReyes2023}. This created a \textbf{substitution effect}: high-performing low-income students now choose between a free elite public seat (Quotas) and a private scholarship (ProUni). This interaction motivates the inclusion of "Public University Expansion" controls in the empirical strategy.

\subsection{Legal Formalia: The Constitutionality of Race}

The shift from a "color-blind" to a "race-conscious" state was solidified by the Supreme Federal Court (\acrshort{stf}).

\subsubsection{ADPF 186 and Material Equality}
The constitutionality of racial quotas was challenged in the case \textbf{ADPF 186}. In a landmark 2012 ruling, the Court voted unanimously to uphold affirmative action. The prevailing opinion explicitly rejected "formal equality" in favor of \textbf{"material equality"}, arguing that the state has a duty of reparative justice to groups historically marginalized by structural racism \cite{STF2012}.

\subsubsection{From Self-Declaration to Heteroidentification}
While ProUni initially relied on self-declaration, reports of fraud and the phenomenon of "money whitening"—where higher-income individuals reclassify themselves to access benefits—undermined the policy's targeting \cite{Telles2004}. This led to the establishment of \textbf{Bancas de Heteroidentificação} (Heteroidentification Commissions) in public universities. These commissions verify phenotype, shifting the institutional definition of ethnicity from subjective identity to administratively verifiable traits, a necessary formalization to enforce the distribution of rights in the "Belíndia" context \cite{RiehlMachadoReyes2023}.